\documentclass{report}
\usepackage{amsmath,amssymb}
\setlength{\parindent}{0mm}
\setlength{\parskip}{1em}
\begin{document}
\begin{center}
\rule{6in}{1pt} \
{\large Rainer Hartmann \\
{\bf An efficient matrix-vector-multiplication for solving high dimensional Schr\"odinger problems using prewavelets and sparse grids}}

Cauerstr 11 \\ D-91058 Erlangen
\\
{\tt rainer.hartmann@fau.de}\\
Christoph Pflaum\end{center}

Solving the Schr\"odinger equation, the basic equation of quantum
mechanics, is an active area of research spanning the fields of
chemistry, physics, and applied mathematics.
The main problem is that this equation is an equation in $ d = 3n$ space
dimensions for a system consisting of $n$ electrons and nuclei.
The so-called ``curse of dimensionality`` prohibits direct approximation
techniques for even reasonably small systems.
A tensor product grid with $ N $ grid points in one dimension consists of
$O \left(N^d\right)$ grid points.
Therefore, it is not possible to solve partial differential equations in
higher dimensions using full grids.
Hence discretization grids are not applied to problems with dimension $d \geq 4$.

Sparse grids can be used to efficiently discretize second order elliptic
differential equations on a $d$-dimensional cube.
The first discretizations of PDEs on sparse grids using the Ritz-Galerkin
approach were restricted to constant coefficients on cubical domains.
Nevertheless, these original discretizations cannot easily be extended to
obtain a computational efficient discretization of partial differential
equations with variable coefficients.
An extension to variable coefficients is essential for applying sparse
grids to PDEs used for problems in the field of Physics and Engineering.

A new sparse grid discretization for Helmholtz equations with a variable
coefficient has been developed.
To reduce the complexity of the sparse grid discretization matrix, we
apply prewavelets and the semi-orthogonality property.
Semi-orthogonality can be treated as a special kind of matrix compression.
An efficient algorithm for matrix vector multiplication using a
Ritz-Galerkin discretization is presented.
This algorithm is based on standard 1-dimensional restrictions and
prolongations, a simple prewavelet stencil, and the classical operator
dependent stencil for multilinear finite elements.
Simulation results show a convergence of the discretization according to
the approximation properties of the finite element space.
The linear equation system is solved by a preconditioned conjugate gradient method.
The condition number of the stiffness matrix can be bounded below $10$
using a standard diagonal preconditioner for a three dimensional test
problem.
Numerical simulation results are presented for a 3-dimensional problem on
a curvilinear bounded domain and for a 6-dimensional problem with
variable coefficients.
Furthermore, simulation results for homogeneous and inhomogeneous
boundary conditions are presented for block-structured grids.


\end{document}
